\chapter{Bézout, Gauss et équations diophantiennes}
\section*{Exercice 30.1 -- Bézout explicite}
Euclide : $35=22+13$, $22=13+9$, $13=9+4$, $9=2\times4+1$. Remonter donne $1=9-2(13-9)=3(22-13)-2\times13=3\times22-5(35-22)=8\times22-5\times35$. Ainsi $u=-5,v=8$.

\section*{Exercice 30.2 -- Inverse modulaire}
$7\times8=56\equiv1\pmod{11}$, donc multiplier par $8$ donne $x\equiv24\equiv2\pmod{11}$. Toutes les solutions sont $x=2+11k$, $k\in\Z$.

\section*{Exercice 30.3 -- Équation diophantienne}
Diviser par le pgcd $3$ donne $2x+3y=5$. Une solution est $(1,1)$. Toute autre vérifie $2(x-1)=-3(y-1)$ ; Gauss impose $3\mid(x-1)$. Donc $x=1+3k$, $y=1-2k$, $k\in\Z$. La substitution prouve réciproquement que tous ces couples conviennent.

\section*{Exercice 30.4 -- Aucune solution}
Le membre de gauche est divisible par $7$, mais $5$ ne l'est pas. Il n'existe donc aucune solution. Le critère général exige que le pgcd des coefficients divise le second membre.

\section*{Exercice 30.5 -- Gauss et puissances}
Bézout donne $ua+vb=1$. Au carré, $1=a(u^2a+2uvb)+v^2b^2$. Cette identité de Bézout pour $a,b^2$ prouve qu'ils sont premiers entre eux. Autre voie : aucun nombre premier ne peut diviser simultanément $a$ et $b^2$.

\section*{Exercice 30.6 -- Deux congruences}
Écrire $x=2+3k$. Modulo $5$, $3k\equiv1$, donc $k\equiv2$ car l'inverse de $3$ est $2$. Ainsi $x=8+15t$, $t\in\Z$. La substitution vérifie les deux conditions. Le module commun est le produit car $3$ et $5$ sont premiers entre eux.
